\section{The language surface and the register-level IR}
\label{sec:language}

This section presents \pqlang{} from the bottom of the stack upward: the
embedding model and register views (\secref{sec:builder}), instruction and
control structure (\secref{sec:instructions}), the register-level IR and its
invariants (\secref{sec:rir}), capability checking
(\secref{sec:capabilities}), workspace management (\secref{sec:layout}), and
the backend-facing architecture (\secref{sec:architecture}, \figref{fig:architecture}).

\subsection{An embedded builder over typed register views}
\label{sec:builder}

\pqlang{} is embedded in Python (version 3.11 or later) and has no
third-party runtime dependencies: the core is a builder library whose output
is data, not an executed program.  A module is built by declaring typed
quantum registers and emitting operations on \emph{views} of them:

\begin{lstlisting}[style=pyqpython, caption={A minimal \pqlang{} program (from the doctested tutorial suite).  Register index~0 is the least significant bit; \texttt{fuse} concatenates views across root registers; \texttt{simulate} is the dependency-free reference executor.}, label={lst:bell}]
from pyqecclang import Bits, Builder, export_originir, simulate

b = Builder("bell_pair", {"pair": Bits(2)})
b.h(b["pair"][0])
b.xor(b["pair"][0], b["pair"][1])
program = b.finish().program()

print(simulate(program).amplitudes)      # {(0, 0): 1, (3, 3): 1} / sqrt 2
print(export_originir(program).text)
\end{lstlisting}

Registers carry one of four storage kinds---\texttt{bits}, \texttt{uint},
\texttt{sint}, or \texttt{rational}, each with a width in $[0,64]$---chosen
so that downstream register-level backends can map them to native integer
registers; total program width is unbounded.  A \emph{view} (an internal
\texttt{Ref}) is a list of \emph{spans} $(\text{register}, \text{start},
\text{width})$: slices of root registers and \texttt{fuse}d concatenations
across them, possibly \texttt{reinterpret}ed to a different storage kind.
Views are formation-time objects: building, slicing, or reinterpreting a view
emits \emph{no} quantum instructions.  Consequently the IR never contains
physical qubit indices at all; register names, widths, and views survive
verbatim to backends (R3).

Builders enforce a strict \emph{alias-freedom ABI}: operands of a single
instruction must not overlap, control registers are protected inside
controlled blocks, and validation rejects violations before any backend sees
the program.  Aliasing analyses of this kind are a prerequisite for the
uncompute-and-prepare discipline that scientific-computing circuits use
pervasively.

\subsection{Instructions and control flow}
\label{sec:instructions}

The closed instruction set is deliberately small: broadcast unary gates
(\texttt{h}, \texttt{x}, \texttt{y}, \texttt{z}, \texttt{s}, \texttt{t},
\texttt{rx}, \texttt{ry}, \texttt{rz}, \texttt{phase},
\texttt{global\_phase}), reversible bit operations (\texttt{xor}, i.e., a
CNOT per bit, and \texttt{swap}), constant modular addition
(\texttt{add\_const}), and QRAM \texttt{Load}.  All richer
structure---arithmetic, oracles, entire solver kernels---is expressed as
\emph{module calls}, so that the IR's semantics remain auditable at a fixed
size.

Control and repetition are structured blocks, entered as Python context
managers:

\begin{lstlisting}[style=pyqpython, caption={Symbolic scale: a legitimate program with a repetition count of $2^{40}$.  The finished program contains a single \texttt{Repeat} node; export remains logarithmic (this is asserted in the test suite).}, label={lst:repeat}]
b = Builder("main", {"word": UInt(8)})
with b.repeat(1 << 40):
    b.add_const(b["word"], 1)
p = b.finish().program()

assert len(p.main.body) == 1                    # one Repeat node
assert export_originir(p).text.count("DEF ") < 45   # log-time export
\end{lstlisting}

Repetition counts may range up to $2^{63}-1$; the body must be structurally
valid even at count zero; and---as a normative rule of the IR
specification---generation and serialization must never replicate a repeated
body according to its count (R5).  Guards against runaway expansion are cost
models (\texttt{expanded\_steps} budgets consumed by the lazy expansion
iterator), not eager materialization.  \texttt{Control} blocks are coherent:
control values are compile-time constants attached to protected registers,
and nested controls must not overlap.  \texttt{Adjoint} blocks are symbolic:
their semantics are defined without rewriting the body, so the adjoint of a
composition is available without duplicating it.

\begin{figure}[t]
  \centering
  \begin{quantikz}[remember picture, every matrix/.append style={name=pqRep}]
  \lstick{$\mathtt{q}_{0}$} & \gate{H} & \ctrl{1} \\
  \lstick{$\mathtt{q}_{1}$} & \qw & \targ{} \\
  \end{quantikz}
  \begin{tikzpicture}[remember picture, overlay]
  \node[draw=black!55, dashed, rounded corners=2pt, inner sep=3pt, fit=(pqRep-1-1)(pqRep-1-2)(pqRep-2-2), label={[font=\footnotesize\color{black!65}]above:$\times 2^{40}$}] {};
  \end{tikzpicture}
  \caption{A \texttt{Repeat} block drawn by the quantikz export.  The dashed
  group is a single \texttt{Repeat} node with a static count of $2^{40}$; the
  body appears exactly once in the program, its serialization, and every
  export---expansion, if a backend performs it, is logarithmic and leaves no
  trace in the IR (\lstref{lst:repeat}).}
  \label{fig:repeat-circuit}
\end{figure}

\subsection{The register-level IR (RIR)}
\label{sec:rir}

The IR is the architectural center of the language.  A program is
$\langle\text{entry}, \text{modules}, \text{version}\rangle$; a module is a
named, typed interface (ordered quantum registers, ordered QRAM resources,
optional private \emph{locals}) plus a body---an immutable tuple of
instructions---or \emph{no body at all}, in which case the module is an
\emph{open declaration}: an oracle slot (\secref{sec:oracles}).
Table~\ref{tab:rir-nodes} summarizes the node kinds.

\begin{table}[t]
  \centering
  \caption{RIR 0.3 objects and instruction nodes.  All nodes are immutable;
  bodies are tuples.  \texttt{Module} with \texttt{body} $=\varnothing$ is an
  open (oracle) declaration.}
  \label{tab:rir-nodes}
  \small
  \begin{tabular}{@{}llp{9.6cm}@{}}
    \toprule
    \tabhead{Object} & \tabhead{Kind} & \tabhead{Contents / semantics} \\
    \midrule
    \texttt{Program} & container & entry module, module table, IR version \\
    \texttt{Module} & definition & name; registers (quantum interface); resources (QRAM); \texttt{locals} (zero-in/zero-out workspace); attributes (e.g.\ \texttt{be\_alpha}, paradigm, provenance); body or $\varnothing$ \\
    \texttt{RegType} & type & kind (\texttt{bits}/\texttt{uint}/\texttt{sint}/\texttt{rational}) and width $0..64$; total width unbounded \\
    \texttt{Ref} & view & tuple of spans $(\text{register}, \text{start}, \text{width})$; declarative, never contiguous physical qubits \\
    \texttt{QRAM} & resource & address width, data width; runtime data binding \\
    \midrule
    \texttt{Primitive} & instruction & broadcast unary gate, \texttt{xor}, \texttt{swap}, \texttt{add\_const}, \texttt{global\_phase} \\
    \texttt{Load} & instruction & QRAM access $|a\rangle|d\rangle \mapsto |a\rangle|d \oplus M[a]\rangle$; self-inverse \\
    \texttt{Call} & instruction & module invocation; kind-and-width-exact, alias-free argument matching; modules are shared, never copied \\
    \texttt{Repeat} & instruction & static count $0..2^{63}{-}1$; body preserved, never replicated \\
    \texttt{Control} & instruction & coherent conditioning on protected registers with compile-time values \\
    \texttt{Adjoint} & instruction & symbolic inversion; no body rewriting \\
    \bottomrule
  \end{tabular}
\end{table}

Four invariants define the IR's contract with the rest of the system:

\begin{enumerate}
\item \textbf{Modularity end-to-end (R3, R5).}  \texttt{Call}, \texttt{Repeat},
  \texttt{Control}, and \texttt{Adjoint} are preserved in the IR, in its JSON
  serialization, and in exports.  Lowering rules permit backends to expand
  register operations to physical gates, but module definitions and calls
  must be preserved, and any expansion happens in the backend without
  feeding back into the RIR.
\item \textbf{Purity.}  RIR represents quantum operations after compile-time
  parameter evaluation---all widths, integers, and angles are concrete.  It
  contains no measurement, reset, runtime classical feedback, or opaque host
  callbacks (R7); readout is planned host-side
  (\secref{sec:readout}).
\item \textbf{Canonical, hardened serialization.}  Programs serialize to
  canonical JSON (sorted keys, deterministic output) under a versioned
  JSON~Schema (currently RIR~0.3, Draft~2020-12).  Decoding is strict:
  unknown tags or fields, duplicate keys, non-finite floats, and unknown
  versions are rejected.  Round-trips are byte-identical, and open modules
  round-trip unchanged.  Readers accept prior format versions, so stored
  programs are forward-compatible artifacts of record.
\item \textbf{Attributes never change instruction semantics.}  Domain
  metadata---block-encoding factors, paradigm tags, implementation-status and
  provenance strings (R6)---ride along as attributes without affecting the
  operational meaning of the module they annotate.
\end{enumerate}

Export to the modular textual backend (OriginIR-ext) illustrates the
preservation contract concretely.  \lstref{lst:originir} shows the exported
text of a QRAM-lookup demo program: module definitions (\texttt{DEF}),
calls, and the \texttt{QRAMDECL} resource declaration survive as first-class
textual constructs; nothing is inlined.

\begin{lstlisting}[style=pyqir, caption={OriginIR-ext export of a nested QRAM-lookup program (verbatim from the repository's example artifacts).  \texttt{DEF} blocks, module calls, and the \texttt{QRAMDECL} resource are preserved; no call is inlined.  Content hashes in definition names make exports deterministic and diff-stable.}, label={lst:originir}]
QRAMDECL ram_values 2,3
QINIT 5
CREG 0
DEF m_lookup_2_3_54524a0d1220cf4bd723f2aa(v_address[2], v_data[3])
ram_values v_address[0], v_address[1], v_data[0], v_data[1], v_data[2]
ENDDEF
DEF m_wrapped_lookup_16dc443ca6e4d78bd9ef03a6(v_a[2], v_d[3])
m_lookup_2_3_54524a0d1220cf4bd723f2aa(v_a[0], v_a[1], v_d[0], v_d[1], v_d[2])
ENDDEF
DEF m_demo_b89906731f4fe79582ed7931(v_address[2], v_data[3])
H v_address[0]
H v_address[1]
X v_data[0]
m_wrapped_lookup_16dc443ca6e4d78bd9ef03a6(v_address[0], v_address[1],
  v_data[0], v_data[1], v_data[2])
ENDDEF
m_demo_b89906731f4fe79582ed7931(q[0], q[1], q[2], q[3], q[4])
\end{lstlisting}

\begin{figure}[t]
  \centering
  \footnotesize
  \begin{quantikz}[remember picture, every matrix/.append style={name=pqDemo}]
  \lstick{$\mathtt{address}_{0}$} & \gate{H} & \qw & \gate[wires=5]{\mathtt{wrapped\_lookup}} \\
  \lstick{$\mathtt{address}_{1}$} & \gate{H} & \qw & \qw \\
  \lstick{$\mathtt{data}_{0}$} & \qw & \gate{X} & \qw \\
  \lstick{$\mathtt{data}_{1}$} & \qw & \qw & \qw \\
  \lstick{$\mathtt{data}_{2}$} & \qw & \qw & \qw \\
  \end{quantikz}\par\vspace{0.3em}
  $\downarrow$\par\vspace{0.3em}
  \begin{quantikz}[remember picture, every matrix/.append style={name=pqWrap}]
  \lstick{$\mathtt{a}_{0}$} & \gate[wires=5]{\mathtt{lookup\_2\_3}} \\
  \lstick{$\mathtt{a}_{1}$} & \qw \\
  \lstick{$\mathtt{d}_{0}$} & \qw \\
  \lstick{$\mathtt{d}_{1}$} & \qw \\
  \lstick{$\mathtt{d}_{2}$} & \qw \\
  \end{quantikz}\par\vspace{0.3em}
  $\downarrow$\par\vspace{0.3em}
  \begin{quantikz}[remember picture, every matrix/.append style={name=pqLookup}]
  \lstick{$\mathtt{address}_{0}$} & \gate[wires=5]{\mathrm{QRAM}_{\mathtt{table}}} \\
  \lstick{$\mathtt{address}_{1}$} & \qw \\
  \lstick{$\mathtt{data}_{0}$} & \qw \\
  \lstick{$\mathtt{data}_{1}$} & \qw \\
  \lstick{$\mathtt{data}_{2}$} & \qw \\
  \end{quantikz}
  \caption{The three modules of \lstref{lst:originir}, drawn independently by
  the quantikz export of the same program: \texttt{demo} (top) prepares
  address and data and calls \texttt{wrapped\_lookup} (middle), which calls
  \texttt{lookup\_2\_3} (bottom), the module that owns the \texttt{Load} on
  the \texttt{table} resource.  At every level a call is a box with a name
  and a typed interface; nothing is expanded, matching the \texttt{DEF} chain
  of the textual export.}
  \label{fig:nested-circuit}
\end{figure}

\subsection{Adjoint and control as a capability discipline}
\label{sec:capabilities}

Whether an object may be used as $U^{\dagger}$ or as controlled-$U$ is
inferred bottom-up from its parts: a primitive gate admits both; a
\texttt{Call} admits adjoint (respectively, control) only if its callee does;
open oracle declarations advertise their capabilities as part of the
declaration, and the binder checks that the provided implementation actually
possesses them (\secref{sec:oracles}).  Validation then checks every
contextual use: wrapping a call in an \texttt{Adjoint} block demands
invertibility of the callee \emph{at composition time}.  This is a modal
capability system in the practical sense---R4---and it eliminates a common
class of late-stage backend failures.

\subsection{Workspaces and qubit economy}
\label{sec:layout}

Modules may declare private \texttt{locals}: workspace registers that must be
returned to zero at exit (enforced during execution; dirty workspaces raise).
A layout pass computes sequential-reuse qubit budgets for module-private
locals, so that deeply composed programs remain qubit-frugal without the
programmer managing temporary registers by hand.  The math-function compiler
(\secref{sec:mathfunc}) leans on this discipline heavily: its compute /
XOR-update / uncompute pattern is only sound because workspace zeroing is an
enforced ABI, not a convention.

\subsection{Architecture: generation, storage, export, execution}
\label{sec:architecture}

\figref{fig:architecture} shows the complete architecture.  The separation
follows the three-phase paradigm of \secref{sec:requirements}: Python
generation (builders, algorithm libraries, protocols and contracts), RIR
storage (validation, linking, serialization, layout), and backend export or
execution.  Native backends are imported lazily at execution entry
only---generating and serializing RIR requires no simulator installation, and
the core has no backend dependency.  Backends currently comprise:

\begin{itemize}
\item \texttt{export\_originir}: modular textual export (originally to the
  OriginIR-ext dialect understood by the \emph{uniqc} simulator), with an
  additional strict lowering of the whole artifact to the
  \{\textsc{toffoli}, $U_3$, \textsc{cz}\} gate set with ancilla pooling;
\item \texttt{simulate}: a dependency-free sparse-statevector reference
  executor with step and state budgets, enforcing workspace-zeroing and
  memory-binding rules;
\item \texttt{run\_pysparq}: register-level native execution on PySparQ,
  including a registry for \emph{native operators}---module-level
  implementations (e.g.\ large reversible arithmetic) compiled on demand to
  C++ and charged unit cost rather than expanded;
\item \texttt{quantikz}: circuit-diagram export for documentation and
  review, preserving the same modularity discipline (calls as named boxes,
  open slots dashed, structured blocks grouped); every circuit figure in
  this paper is generated by it.
\end{itemize}

Measurement, reset, and statistics are planned as host-side
\emph{readout actions} appended after export (\secref{sec:readout}).

\begin{figure}[t]
  \centering
  \begin{tikzpicture}[
      font=\small,
      box/.style={draw=black!60, rounded corners=2pt, fill=black!3, align=center, inner sep=4pt},
      phase/.style={font=\footnotesize\bfseries, text=black!45, rotate=90},
      sepline/.style={black!22, dotted},
      arr/.style={-{Stealth[length=2.2mm]}, black!60},
      darr/.style={-{Stealth[length=2.2mm]}, black!35, dashed},
    ]
    % Phase tags and separators
    \node[phase] at (-6.85, 2.7) {generate};
    \node[phase] at (-6.85, -1.6) {store \& link};
    \node[phase] at (-6.85, -5.45) {export / execute};
    \draw[sepline] (-6.4, -0.5) -- (7.15, -0.5);
    \draw[sepline] (-6.4, -4.05) -- (7.15, -4.05);

    % Generation phase
    \node[box, fill=black!8] (python) at (0, 3.5) {\textbf{Python host}};
    \node[box] (builder) at (0, 2.1)
      {\textbf{Builder}\\\footnotesize typed registers, views\\\footnotesize control / repeat / adjoint};
    \node[box] (libs) at (-4.5, 2.1)
      {\textbf{algorithm libraries}\\\footnotesize protocols, contracts\\\footnotesize open oracle declarations};
    \draw[arr] (python) -- (builder);
    \draw[arr] (libs) -- (builder);

    % Storage phase
    \node[box, fill=black!8] (rir) at (0, -1.6)
      {\textbf{RIR 0.3}\\\footnotesize modules (open \& closed)\\\footnotesize calls / repeats symbolic};
    \node[box] (validate) at (-4.5, -1.6)
      {\footnotesize validation\\\footnotesize (aliases, capabilities)};
    \node[box] (bindbox) at (4.5, -1.6)
      {\footnotesize \texttt{bind} linking\\\footnotesize (oracles, QRAM hoisting)};
    \node[box] (json) at (0, -3.1)
      {\footnotesize canonical JSON + schema\\\footnotesize (deterministic, hardened)};
    \draw[arr] (builder) -- (rir);
    \draw[arr] (validate) -- (rir);
    \draw[arr] (rir) -- (bindbox);
    \draw[arr] (bindbox.south) to[bend left=25] (rir.south);
    \draw[arr] (rir) -- (json);

    % Backends
    \node[box] (originir) at (-5.55, -5.4)
      {\footnotesize \textbf{OriginIR-ext}\\\scriptsize modular text\\\scriptsize (+ Toffoli/$U_3$/CZ)};
    \node[box] (sim) at (-1.9, -5.4) {\footnotesize \textbf{reference}\\\scriptsize \texttt{simulate}};
    \node[box] (pysparq) at (1.9, -5.4)
      {\footnotesize \textbf{PySparQ native}\\\scriptsize registers + C++ ops};
    \node[box, fill=black!8] (readout) at (5.55, -5.4)
      {\footnotesize \textbf{host readout}\\\scriptsize measure / stats};
    \draw[arr] (json.south) to[bend right=10] (originir.north);
    \draw[arr] (json.south) -- (sim.north);
    \draw[arr] (json.south) to[bend left=10] (pysparq.north);
    \draw[darr] (pysparq.east) -- (readout.west);
    \draw[darr] (sim.south) -- ++(0, -0.55) -| (readout.south);
  \end{tikzpicture}
  \caption{\pqlang{} architecture.  Programs are generated by the embedded
  builder and algorithm libraries, stored as RIR (open oracle slots included),
  optionally linked by \texttt{bind}, serialized to canonical schema-validated
  JSON, and only then exported or executed.  Native backends are imported
  lazily at execution entry; the core has no backend dependency.  Readout is
  a host-side plan, not an IR construct.}
  \label{fig:architecture}
\end{figure}
